\section{线性分类器：深度学习的基石}

线性分类器是深度学习中\textbf{最重要的基础模块}。几乎所有的神经网络架构都以线性变换为核心组件。

\subsection{从非参数到参数化方法}

KNN 是\textbf{非参数方法}——它不学习任何参数，而是直接存储所有训练数据。线性分类器是\textbf{参数化方法}——它学习一组参数 $W$（权重矩阵），将输入图像映射到输出类别分数。

\begin{figure}[H]
\centering
\includegraphics[width=0.85\textwidth]{slides-images/slide-050.jpg}
\caption{参数化方法：线性分类器 $f(x, W) = Wx$，输入图像 $x$ 通过权重矩阵 $W$ 映射到类别分数\protect\footnotemark}
\end{figure}
\footnotetext{Slides 第53页。}

\begin{importantbox}{参数化方法 vs 非参数方法}
\begin{itemize}
    \item \textbf{KNN（非参数方法）}：训练 $O(1)$，预测 $O(N)$。模型就是全部训练数据。
    \item \textbf{线性分类器（参数化方法）}：训练时间较长（需要优化），但预测极快——只需一次矩阵乘法。模型是学习到的参数 $W$ 和 $b$，训练数据在预测时不再需要。
\end{itemize}
\end{importantbox}

\subsection{线性分类器的数学定义}

\begin{figure}[H]
\centering
\includegraphics[width=0.95\textwidth]{slides-images/slide-053.jpg}
\caption{线性分类器示意：将 $32 \times 32 \times 3 = 3072$ 维的输入向量通过 $W$ 映射到 10 个类别分数\protect\footnotemark}
\end{figure}
\footnotetext{Slides 第55页。}

线性分类器的核心公式为：

$$f(x, W) = Wx + b$$

其中：
\begin{itemize}
    \item $x \in \mathbb{R}^{D \times 1}$：输入图像展平为列向量（例如 CIFAR-10 中 $D = 32 \times 32 \times 3 = 3072$）
    \item $W \in \mathbb{R}^{C \times D}$：权重矩阵（$C$ 为类别数，例如 $C = 10$）
    \item $b \in \mathbb{R}^{C \times 1}$：偏置向量，与输入无关
    \item $f(x, W) \in \mathbb{R}^{C \times 1}$：每个类别的分数
\end{itemize}

\begin{figure}[H]
\centering
\includegraphics[width=0.95\textwidth]{slides-images/slide-056.jpg}
\caption{线性分类器的代数计算过程：以 $2 \times 2$ 图像（4个像素）和 3 个类别（cat, dog, ship）为例\protect\footnotemark}
\end{figure}
\footnotetext{Slides 第58页。}

\begin{knowledgebox}{偏置项 $b$ 的作用}
偏置向量 $b$ 提供了一个\textbf{与输入无关}的类别偏好。例如，如果训练数据中猫的样本远多于其他类别，偏置项会自动调整，使``猫''类的分数整体偏高。在几何上，偏置项允许决策超平面不必经过原点，从而提供更灵活的分类边界。
\end{knowledgebox}

\subsection{线性分类器的三种解读视角}

理解线性分类器有三种互补的视角：代数视角、视觉视角和几何视角。

\subsubsection{代数视角（Algebraic Viewpoint）}

\begin{figure}[H]
\centering
\includegraphics[width=0.95\textwidth]{slides-images/slide-056.jpg}
\caption{代数视角：矩阵乘法加偏置，将输入向量线性映射到类别分数向量\protect\footnotemark}
\end{figure}
\footnotetext{Slides 第58页。}

从代数角度看，$f = Wx + b$ 就是一个简单的矩阵-向量乘法加上偏置。$W$ 的第 $i$ 行 $w_i$ 与输入 $x$ 做内积再加上 $b_i$，得到第 $i$ 类的分数。

\subsubsection{视觉视角（Visual Viewpoint）}

\begin{figure}[H]
\centering
\includegraphics[width=0.95\textwidth]{slides-images/slide-060.jpg}
\caption{视觉视角：$W$ 的每一行可以还原为一个``模板图像''，代表该类别的典型外观\protect\footnotemark}
\end{figure}
\footnotetext{Slides 第62页。}

$W$ 的每一行 $w_i \in \mathbb{R}^D$ 可以被 reshape 为与输入图像相同的形状（如 $32 \times 32 \times 3$），得到一个\textbf{模板图像}（template）。分类过程本质上是用每个类别的模板与输入图像做\textbf{匹配}（内积越大，匹配度越高）。

在 CIFAR-10 上训练后，可以可视化这些模板：
\begin{itemize}
    \item ``car'' 模板呈现出红色的车辆正面轮廓
    \item ``airplane'' 模板呈现出蓝色背景上的飞机形状
    \item ``horse'' 模板呈现出绿色背景上的马的轮廓
\end{itemize}

\begin{warningbox}{线性分类器只能学习一个模板}
线性分类器的根本局限在于：每个类别只有\textbf{一个}模板（$W$ 的一行）。如果同一类别有多种外观（如红色汽车和蓝色汽车），线性分类器只能学习到一个``平均''的模板。这就是为什么我们后续需要神经网络——它们可以学习多层次、多模式的特征表示。
\end{warningbox}

\subsubsection{几何视角（Geometric Viewpoint）}

\begin{figure}[H]
\centering
\includegraphics[width=0.95\textwidth]{slides-images/slide-060.jpg}
\caption{几何视角：线性分类器在像素空间中为每对类别画一条线性决策边界\protect\footnotemark}
\end{figure}
\footnotetext{Slides 第63页。}

在像素空间中，线性分类器对每个类别学习一个\textbf{线性决策边界}（在 2D 中是直线，在高维中是超平面）。$W$ 的第 $i$ 行定义了超平面的法向量，$b_i$ 定义了超平面的偏移量。

\begin{figure}[H]
\centering
\includegraphics[width=0.95\textwidth]{slides-images/slide-063.jpg}
\caption{几何视角的高维表示：每个分类器定义一个超平面，将空间分为正（属于该类）和负（不属于该类）两侧\protect\footnotemark}
\end{figure}
\footnotetext{Slides 第65页。}

\subsection{线性分类器的局限性}

\begin{figure}[H]
\centering
\includegraphics[width=0.95\textwidth]{slides-images/slide-064.jpg}
\caption{线性分类器无法处理的三种情况：XOR 问题、环形分布、多模态分布\protect\footnotemark}
\end{figure}
\footnotetext{Slides 第66页。}

存在多种线性分类器\textbf{无法处理}的数据分布模式：

\begin{enumerate}
    \item \textbf{XOR 问题}：两个类别分别占据第一三象限和第二四象限，无法用一条直线分开
    \item \textbf{环形分布}：一个类别包围另一个类别
    \item \textbf{多模态分布}：一个类别的数据分布在空间中多个不相连的区域
\end{enumerate}

\begin{warningbox}{线性不可分问题}
当数据不是\textbf{线性可分}的时候，线性分类器无论怎么优化参数都无法得到正确的分类结果。解决方法包括：（1）引入非线性特征变换；（2）堆叠多个线性层并加入非线性激活函数——这就是\textbf{神经网络}的基本思想。
\end{warningbox}

\subsection{本章小结}

线性分类器通过学习权重矩阵 $W$ 和偏置 $b$，实现从输入图像到类别分数的参数化映射。它有三种等价的理解方式：代数上是矩阵乘法；视觉上是模板匹配；几何上是超平面分割。线性分类器训练后只保留参数，预测极快，但其表达能力有限——只能学习线性决策边界，每类仅一个模板。它是构建更复杂神经网络的基石。

%% ========== Part 5: 损失函数 ==========

